\documentclass[12pt,a4paper]{ctexart}

% ==== 字体与编码 ====
\usepackage{fontspec}
\setmainfont{Times New Roman}
\setsansfont{Arial}
\setmonofont{Consolas}
\setCJKmainfont{SimSun}[AutoFakeBold=2.5]
\setCJKsansfont{SimHei}
\setCJKmonofont{KaiTi}

% ==== 数学 ====
\usepackage{amsmath,amssymb,amsthm}
\usepackage{mathtools}

% ==== 页面布局 ====
\usepackage[top=2.5cm,bottom=2.5cm,left=2.8cm,right=2.8cm]{geometry}

% ==== 超链接 ====
\usepackage[colorlinks=true,linkcolor=blue,citecolor=blue,urlcolor=blue]{hyperref}

% ==== 交换图 ====
\usepackage{tikz-cd}

% ==== 列表 ====
\usepackage{enumitem}

% ==== 表格 ====
\usepackage{booktabs,array}

% ==== 定理环境 ====
\theoremstyle{definition}
\newtheorem{definition}{定义}[section]
\newtheorem{example}[definition]{例}
\newtheorem{remark}[definition]{注记}
\newtheorem{theorem}[definition]{定理}
\newtheorem{proposition}[definition]{命题}
\newtheorem{corollary}[definition]{推论}

% ==== 自定义命令 ====
\newcommand{\cat}[1]{\mathbf{#1}}
\newcommand{\Set}{\cat{Set}}
\newcommand{\id}{\mathrm{id}}
\newcommand{\ob}[1]{\mathrm{Ob}(#1)}
\newcommand{\Hom}{\mathrm{Hom}}
\newcommand{\dashvadj}{\dashv}

\title{\textbf{伴随函子的双重视角与单子/余单子的自然诱导}\\[0.3em]
\large ——从单位--余单位到单子公理的完整推演}
\author{范畴论学习笔记 · 范畴专栏}
\date{}

\begin{document}
\maketitle

\tableofcontents
\newpage

% ==================================================================
\section{引言：一个概念，两种定义}
% ==================================================================

上一讲我们给出了伴随函子的严格定义——基于 Hom-集的双函子自然同构
\[
\Phi_{X,Y}: \Hom_{\cat{D}}(F(X), Y) \cong \Hom_{\cat{C}}(X, G(Y))
\]
这一定义。然而，范畴论中还存在另一个等价的、同样核心的定义：\textbf{单位与余单位}视角。

这两种定义描述的是同一个数学本质，它们严格等价，且各有优势：
\begin{itemize}[nosep]
    \item \textbf{Hom-集版本}：直观反映“态射集的一一对应”，适合入门和具体计算；
    \item \textbf{单位--余单位版本}：从全局自然变换的角度定义，适合证明普遍性质和构造单子。
\end{itemize}

更重要地，\textbf{任意一对伴随函子 $(F, G, \eta, \varepsilon)$ 都自然诱导出范畴 $\cat{C}$ 上的一个单子（Monad）和范畴 $\cat{D}$ 上的一个余单子（Comonad）}。这就是单子和余单子最本源的动机——它们并非人为凭空定义的抽象代数结构，而是伴随函子复合的自然产物。

本文目标：
\begin{enumerate}[nosep]
    \item 严格证明 Hom-集定义与单位--余单位定义的等价性；
    \item 从伴随出发，构造单子 $T = GF$，逐条验证单子公理；
    \item 构造余单子 $S = FG$，验证余单子公理；
    \item 用交换图展示完整的代数结构。
\end{enumerate}

\newpage

% ==================================================================
\section{伴随函子的两个等价定义}
% ==================================================================

\subsection{定义一：Hom-集自然同构（回顾）}

\begin{definition}[Hom-集版本]\label{def:adj-hom}
给定函子 $F: \cat{C} \to \cat{D}$ 和 $G: \cat{D} \to \cat{C}$。$F \dashvadj G$ 当且仅当存在自然同构
\[
\Phi: \Hom_{\cat{D}}(F(-), -) \;\cong\; \Hom_{\cat{C}}(-, G(-))
\]
作为函子范畴 $[\cat{C}^{\mathrm{op}} \times \cat{D}, \Set]$ 中的同构。
\end{definition}

\subsection{定义二：单位与余单位}

\begin{definition}[单位--余单位版本]\label{def:adj-unit}
给定函子 $F: \cat{C} \to \cat{D}$ 和 $G: \cat{D} \to \cat{C}$。$F \dashvadj G$ 当且仅当存在两个自然变换：
\begin{align}
\eta &: \id_{\cat{C}} \Longrightarrow GF \tag{单位/Unit}\\
\varepsilon &: FG \Longrightarrow \id_{\cat{D}} \tag{余单位/Counit}
\end{align}
满足以下\textbf{三角形恒等式（Triangle Identities）}：
\begin{align}
\varepsilon_{F(X)} \circ F(\eta_X) &= \id_{F(X)} \label{eq:tri1-comp}\\
G(\varepsilon_Y) \circ \eta_{G(Y)} &= \id_{G(Y)} \label{eq:tri2-comp}
\end{align}
\end{definition}

\subsection{等价性证明}

\begin{theorem}[两种定义的等价性]\label{thm:equiv}
定义 \ref{def:adj-hom} 与定义 \ref{def:adj-unit} 等价。
\end{theorem}

\begin{proof}
分两个方向证明。

\medskip\noindent\textbf{方向一：从 Hom-集双射 $\Phi$ 构造 $\eta$ 和 $\varepsilon$。}

对任意 $X \in \cat{C}$，取 $Y = F(X)$，考虑双射
\[
\Phi_{X, F(X)}: \Hom_{\cat{D}}(F(X), F(X)) \xrightarrow{\sim} \Hom_{\cat{C}}(X, GF(X))
\]
取恒等态射 $\id_{F(X)}$ 在左边的值，定义
\begin{equation}\label{eq:eta-def}
\eta_X := \Phi_{X, F(X)}(\id_{F(X)}) \; : \; X \to GF(X)
\end{equation}

对任意 $Y \in \cat{D}$，取 $X = G(Y)$，考虑逆映射
\[
\Phi_{G(Y), Y}^{-1}: \Hom_{\cat{C}}(G(Y), G(Y)) \xrightarrow{\sim} \Hom_{\cat{D}}(FG(Y), Y)
\]
取恒等态射 $\id_{G(Y)}$ 在右边的值，定义
\begin{equation}\label{eq:eps-def}
\varepsilon_Y := \Phi_{G(Y), Y}^{-1}(\id_{G(Y)}) \; : \; FG(Y) \to Y
\end{equation}

$\eta$ 和 $\varepsilon$ 的自然性由 $\Phi$ 的自然性推出（见命题 \ref{prop:naturality-eta-eps}）。

验证三角形恒等式：
\[
\varepsilon_{FX} \circ F(\eta_X)
= \Phi_{GFX,FX}^{-1}(\id_{GFX}) \circ F(\Phi_{X,FX}(\id_{FX}))
\]
由统一自然方程 $\Phi_{X',Y'}(g \circ h \circ F(f)) = G(g) \circ \Phi_{X,Y}(h) \circ f$，可证明此式 $= \id_{FX}$（见命题 \ref{prop:triangle-from-Phi}）。

\medskip\noindent\textbf{方向二：从 $\eta, \varepsilon$ 构造 $\Phi$。}

对任意 $f: F(X) \to Y$，定义
\[
\Phi_{X,Y}(f) := G(f) \circ \eta_X \; : \; X \to G(Y)
\]
对任意 $g: X \to G(Y)$，定义
\[
\Phi_{X,Y}^{-1}(g) := \varepsilon_Y \circ F(g) \; : \; F(X) \to Y
\]

验证互为逆映射：
\[
\Phi_{X,Y}^{-1}(\Phi_{X,Y}(f))
= \varepsilon_Y \circ F(G(f) \circ \eta_X)
= f \circ \varepsilon_{FX} \circ F(\eta_X)
= f
\]
\[
\Phi_{X,Y}(\Phi_{X,Y}^{-1}(g))
= G(\varepsilon_Y \circ F(g)) \circ \eta_X
= G(\varepsilon_Y) \circ \eta_{GY} \circ g
= g
\]
$\Phi$ 的自然性由 $\eta, \varepsilon$ 的自然性推出。故 $\Phi$ 是自然同构。
\end{proof}

\begin{proposition}[从 $\Phi$ 构造 $\eta$ 和 $\varepsilon$ 的自然性]\label{prop:naturality-eta-eps}
由 \eqref{eq:eta-def} 和 \eqref{eq:eps-def} 定义的 $\eta$ 和 $\varepsilon$ 是自然变换。
\end{proposition}

\begin{proof}
对任意 $f: X \to X'$，应用 $\Phi$ 的统一自然方程：
\[
\Phi_{X, FX'}(F(f)) = G(F(f)) \circ \eta_X
\]
又 $\Phi_{X, FX'}(F(f)) = \Phi_{X, FX'}(\id_{FX'} \circ \id_{FX} \circ F(f)) = \eta_{X'} \circ f$。
因此 $\eta_{X'} \circ f = GF(f) \circ \eta_X$，$\eta$ 是自然的。$\varepsilon$ 对偶可证。
\end{proof}

\begin{proposition}\label{prop:triangle-from-Phi}
若 $\Phi$ 是自然同构，定义 $\eta$ 和 $\varepsilon$ 如 \eqref{eq:eta-def}、\eqref{eq:eps-def}，则三角形恒等式成立。
\end{proposition}

\begin{proof}
在方向二的构造下，$\Phi$ 和 $\Psi$（由 $\eta, \varepsilon$ 构造）是相同的双射族，而方向二中验证了 $\Phi$ 和 $\Psi$ 互为逆当且仅当三角形恒等式成立。由于 $\Phi$ 是双射，三角形恒等式必然成立。
\end{proof}

综上所述，定理 \ref{thm:equiv} 得证。两种定义完全等价。

\newpage

% ==================================================================
\section{从伴随到单子（Monad）}
% ==================================================================

\subsection{单子的定义（快速回顾）}

\begin{definition}[单子]
范畴 $\cat{C}$ 上的\textbf{单子}由三元组 $(T, \eta, \mu)$ 组成，其中：
\begin{itemize}[nosep]
    \item $T: \cat{C} \to \cat{C}$ 是自函子；
    \item $\eta: \id_{\cat{C}} \Rightarrow T$ 是自然变换（单位）；
    \item $\mu: T^2 \Rightarrow T$ 是自然变换（乘法）；
\end{itemize}
满足以下公理：
\begin{enumerate}[label=(M\arabic*),nosep]
    \item \textbf{结合律：} $\mu \circ (T \mu) = \mu \circ (\mu T)$，即下图交换：
    \[
    \begin{tikzcd}
    T^3 \arrow[r,"T\mu"] \arrow[d,"\mu T"'] & T^2 \arrow[d,"\mu"] \\
    T^2 \arrow[r,"\mu"'] & T
    \end{tikzcd}
    \]
    \item \textbf{单位律：} $\mu \circ (\eta T) = \id_T = \mu \circ (T \eta)$，即下图交换：
    \[
    \begin{tikzcd}
    T \arrow[r,"\eta T"] \arrow[dr,"\id_T"'] & T^2 \arrow[d,"\mu"] & T \arrow[l,"T\eta"'] \arrow[dl,"\id_T"] \\
    & T &
    \end{tikzcd}
    \]
\end{enumerate}
\end{definition}

\subsection{伴随函子自然诱导单子}

\begin{theorem}[伴随 $\Rightarrow$ 单子]\label{thm:adj-monad}
设 $(F, G, \eta, \varepsilon)$ 是伴随函子对 $F \dashvadj G$，其中 $F: \cat{C} \to \cat{D}$，$G: \cat{D} \to \cat{C}$。定义：
\begin{align}
T &:= G \circ F : \cat{C} \to \cat{C} \label{eq:T-def}\\
\eta^{\mathrm{monad}} &:= \eta : \id_{\cat{C}} \Rightarrow T \label{eq:eta-monad}\\
\mu &:= G \varepsilon F : T^2 = GFGF \Longrightarrow GF = T \label{eq:mu-def}
\end{align}
则 $(T, \eta, \mu)$ 是范畴 $\cat{C}$ 上的单子。
\end{theorem}

\begin{remark}[$\mu$ 的展开]
自然变换 $\mu$ 在对象 $X$ 上的分量是：
\[
\mu_X = G(\varepsilon_{F(X)}) : GFGF(X) \to GF(X)
\]
即先应用 $\varepsilon$ 于 $F(X)$：$\varepsilon_{F(X)}: FGF(X) \to F(X)$，再以 $G$ 作用。
\end{remark}

\begin{proof}
验证单子的两条公理。

\medskip\noindent\textbf{（M1）结合律：$\mu \circ (T \mu) = \mu \circ (\mu T)$。}

对任意 $X \in \cat{C}$：
\begin{align*}
\mu_X \circ T(\mu_X) &= G(\varepsilon_{FX}) \circ GFG(\varepsilon_{FX}) \\
&= G\bigl(\varepsilon_{FX} \circ FG(\varepsilon_{FX})\bigr) \\
&= G\bigl(\varepsilon_{FX} \circ \varepsilon_{FGFX}\bigr) \quad (\varepsilon\ \text{的自然性})\\
&= G(\varepsilon_{FX}) \circ G(\varepsilon_{FGFX}) \\
&= \mu_X \circ \mu_{T(X)}
\end{align*}

用交换图表示：
\[
\begin{tikzcd}[column sep=3.5em, row sep=2.8em]
GFGFGFX \arrow[r,"T\mu_X"] \arrow[d,"\mu_{TX}"'] & GFGFX \arrow[d,"\mu_X"] \\
GFGFX \arrow[r,"\mu_X"'] & GFX
\end{tikzcd}
\]

\medskip\noindent\textbf{（M2）单位律：$\mu \circ (\eta T) = \id_T = \mu \circ (T \eta)$。}

对任意 $X \in \cat{C}$：
\begin{align*}
\mu_X \circ \eta_{T(X)} &= G(\varepsilon_{FX}) \circ \eta_{GFX}
= \id_{GFX} \quad (\text{三角形恒等式 \eqref{eq:tri2-comp}},\; Y = FX) \\
\mu_X \circ T(\eta_X) &= G(\varepsilon_{FX}) \circ GF(\eta_X) \\
&= G(\varepsilon_{FX} \circ F(\eta_X))
= G(\id_{FX}) = \id_{GFX} \quad (\text{三角形恒等式 \eqref{eq:tri1-comp}})
\end{align*}

用交换图表示：
\[
\begin{tikzcd}[column sep=3em, row sep=2.2em]
GFX \arrow[r,"\eta_{GFX}"] \arrow[dr,"\id"'] & GFGFX \arrow[d,"G\varepsilon_{FX}"] & GFX \arrow[l,"GF\eta_X"'] \arrow[dl,"\id"] \\
& GFX &
\end{tikzcd}
\]

两条公理验证完毕。故 $(T, \eta, \mu)$ 是单子。
\end{proof}

\subsection{单子公理的宏观交换图}

把上述两条公理以函子级别的 Whiskering 交换图表示：

\[
\begin{tikzcd}[column sep=4.5em, row sep=3em]
T^3 \arrow[r,"T\mu"] \arrow[d,"\mu T"'] & T^2 \arrow[d,"\mu"] \\
T^2 \arrow[r,"\mu"'] & T
\end{tikzcd}
\qquad\Longleftrightarrow\qquad
\begin{tikzcd}[column sep=3.5em, row sep=2.5em]
GFGFGF \arrow[r,"GFG\varepsilon F"] \arrow[d,"G\varepsilon FGF"'] & GFGF \arrow[d,"G\varepsilon F"] \\
GFGF \arrow[r,"G\varepsilon F"'] & GF
\end{tikzcd}
\]

\[
\begin{tikzcd}[column sep=4.5em, row sep=3em]
T \arrow[r,"\eta T"] \arrow[dr,"\id_T"'] & T^2 \arrow[d,"\mu"] & T \arrow[l,"T\eta"'] \arrow[dl,"\id_T"] \\
& T &
\end{tikzcd}
\qquad\Longleftrightarrow\qquad
\begin{tikzcd}[column sep=3em, row sep=2.5em]
GF \arrow[r,"\eta GF"] \arrow[dr,"\id"'] & GFGF \arrow[d,"G\varepsilon F"] & GF \arrow[l,"GF\eta"'] \arrow[dl,"\id"] \\
& GF &
\end{tikzcd}
\]

\begin{remark}[单子即为"由伴随而来的代数"]\label{rmk:monad-origin}
这是范畴论最深刻的事实之一：\textbf{每一个伴随函子对自动产生一个单子}。反过来，\textbf{每一个单子都可以由某个伴随函子对生成}（Kleisli范畴或Eilenberg-Moore代数的构造）。单子并非孤立定义的抽象结构——它的三条公理（结合律、左右单位律）正是伴随的三角形恒等式的直接投影。
\end{remark}

\newpage

% ==================================================================
\section{从伴随到余单子（Comonad）}
% ==================================================================

\subsection{余单子的定义}

余单子是单子的对偶概念（将所有箭头反向）。

\begin{definition}[余单子]
范畴 $\cat{D}$ 上的\textbf{余单子}由三元组 $(S, \varepsilon, \delta)$ 组成，其中：
\begin{itemize}[nosep]
    \item $S: \cat{D} \to \cat{D}$ 是自函子；
    \item $\varepsilon: S \Rightarrow \id_{\cat{D}}$ 是自然变换（余单位）；
    \item $\delta: S \Rightarrow S^2$ 是自然变换（余乘法）；
\end{itemize}
满足以下公理：
\begin{enumerate}[label=(CM\arabic*),nosep]
    \item \textbf{余结合律：} $(S \delta) \circ \delta = (\delta S) \circ \delta$，即下图交换：
    \[
    \begin{tikzcd}
    S \arrow[r,"\delta"] \arrow[d,"\delta"'] & S^2 \arrow[d,"S\delta"] \\
    S^2 \arrow[r,"\delta S"'] & S^3
    \end{tikzcd}
    \]
    \item \textbf{余单位律：} $(\varepsilon S) \circ \delta = \id_S = (S \varepsilon) \circ \delta$，即下图交换：
    \[
    \begin{tikzcd}
    & S \arrow[dl,"\id_S"'] \arrow[d,"\delta"] \arrow[dr,"\id_S"] & \\
    S & S^2 \arrow[l,"\varepsilon S"] \arrow[r,"S\varepsilon"'] & S
    \end{tikzcd}
    \]
\end{enumerate}
\end{definition}

\subsection{伴随函子自然诱导余单子}

\begin{theorem}[伴随 $\Rightarrow$ 余单子]\label{thm:adj-comonad}
设 $(F, G, \eta, \varepsilon)$ 是伴随函子对 $F \dashvadj G$，$F: \cat{C} \to \cat{D}$，$G: \cat{D} \to \cat{C}$。定义：
\begin{align}
S &:= F \circ G : \cat{D} \to \cat{D} \label{eq:S-def}\\
\varepsilon^{\mathrm{comonad}} &:= \varepsilon : S \Rightarrow \id_{\cat{D}} \label{eq:eps-comonad}\\
\delta &:= F \eta G : S = FG \Longrightarrow FGFG = S^2 \label{eq:delta-def}
\end{align}
则 $(S, \varepsilon, \delta)$ 是范畴 $\cat{D}$ 上的余单子。
\end{theorem}

\begin{remark}[$\delta$ 的展开]
自然变换 $\delta$ 在对象 $Y$ 上的分量是：
\[
\delta_Y = F(\eta_{G(Y)}) : FG(Y) \to FGFG(Y)
\]
\end{remark}

\begin{proof}
\textbf{（CM1）余结合律：$(S \delta) \circ \delta = (\delta S) \circ \delta$。}

对任意 $Y \in \cat{D}$：
\begin{align*}
(S\delta)_Y \circ \delta_Y &= FG(F\eta_{GY}) \circ F\eta_{GY}
= F\bigl(GF\eta_{GY} \circ \eta_{GY}\bigr) \\
&= F\bigl(\eta_{GFGY} \circ \eta_{GY}\bigr) \quad (\eta\ \text{的自然性})\\
&= F\eta_{GFGY} \circ F\eta_{GY}
= (\delta S)_Y \circ \delta_Y
\end{align*}

交换图：
\[
\begin{tikzcd}[column sep=4em, row sep=2.8em]
FGY \arrow[r,"\delta_Y"] \arrow[d,"\delta_Y"'] & FGFGY \arrow[d,"S\delta_Y"] \\
FGFGY \arrow[r,"\delta S_Y"'] & FGFGFGY
\end{tikzcd}
\]

\medskip\noindent\textbf{（CM2）余单位律：$(\varepsilon S) \circ \delta = \id_S = (S \varepsilon) \circ \delta$。}

对任意 $Y \in \cat{D}$：
\begin{align*}
(\varepsilon S)_Y \circ \delta_Y &= \varepsilon_{FGY} \circ F\eta_{GY}
= \id_{FGY} \quad (\text{三角形恒等式 \eqref{eq:tri1-comp}},\; X = GY) \\
(S\varepsilon)_Y \circ \delta_Y &= FG(\varepsilon_Y) \circ F\eta_{GY}
= F\bigl(G\varepsilon_Y \circ \eta_{GY}\bigr) \\
&= F(\id_{GY}) = \id_{FGY} \quad (\text{三角形恒等式 \eqref{eq:tri2-comp}})
\end{align*}

交换图：
\[
\begin{tikzcd}
& FGY \arrow[dl,"\id"'] \arrow[d,"\delta_Y"] \arrow[dr,"\id"] & \\
FGY & FGFGY \arrow[l,"\varepsilon_{FGY}"] \arrow[r,"FG\varepsilon_Y"'] & FGY
\end{tikzcd}
\]

故 $(S, \varepsilon, \delta)$ 是余单子。
\end{proof}

\newpage

% ==================================================================
\section{伴随-单子-余单子的统一结构图}
% ==================================================================

\subsection{完整对应关系}

给定伴随 $F \dashvadj G$（$F: \cat{C} \to \cat{D}$），我们获得了：
\begin{itemize}[nosep]
    \item $\cat{C}$ 上的单子 $\mathbb{T} = (GF, \eta, G\varepsilon F)$
    \item $\cat{D}$ 上的余单子 $\mathbb{S} = (FG, \varepsilon, F\eta G)$
\end{itemize}

\[
\begin{tikzcd}[column sep=5em, row sep=4em]
\cat{C} \arrow[r,bend left,"F"] \arrow[loop left,"\mathbb{T}=GF"] & \cat{D} \arrow[l,bend left,"G"] \arrow[loop right,"\mathbb{S}=FG"]
\end{tikzcd}
\]

\begin{table}[h]
\centering
\caption{伴随函子 vs 单子 vs 余单子 结构对应}
\begin{tabular}{@{}l|l|l|l@{}}
\toprule
\textbf{伴随结构} & \textbf{记号} & \textbf{单子 $(GF, \eta, \mu)$} & \textbf{余单子 $(FG, \varepsilon, \delta)$} \\
\midrule
函子对     & $F \dashvadj G$ & $T = GF$ & $S = FG$ \\
单位       & $\eta: \id \Rightarrow GF$ & $\eta$（取用）& $\delta = F\eta G$（余乘法）\\
余单位     & $\varepsilon: FG \Rightarrow \id$ & $\mu = G\varepsilon F$（乘法）& $\varepsilon$（取用）\\
三角形等式 & (\ref{eq:tri1-comp}), (\ref{eq:tri2-comp}) & 单位律 & 余单位律 \\
自然性     & $\eta,\varepsilon$ 的自然性 & 结合律（源于 $\varepsilon$ 自然性） & 余结合律（源于 $\eta$ 自然性） \\
\bottomrule
\end{tabular}
\end{table}

\subsection{公理来源一一对应}

\[
\begin{tikzcd}[column sep=0.8em, row sep=2em]
\text{三角形恒等式} \arrow[dr,Rightarrow] \arrow[ddr,Rightarrow] \\
& \text{左右投射} \arrow[d,mapsto] & \arrow[d,mapsto] \\
& \text{单子单位律} & \text{余单子余单位律} \\[1em]
\varepsilon\ \text{的自然性} \arrow[dr,Rightarrow] \\
& \text{单子结合律} \arrow[d,mapsto] \\
& \mu \circ (T\mu) = \mu \circ (\mu T) \\[1em]
\eta\ \text{的自然性} \arrow[dr,Rightarrow] \\
& \text{余单子余结合律} \arrow[d,mapsto] \\
& (S\delta) \circ \delta = (\delta S) \circ \delta
\end{tikzcd}
\]

% ==================================================================
\section{为什么这是单子的真正动机？}
% ==================================================================

范畴论中，单子（Monad）最初并非凭空定义的一组公理。它的三条公理（结合律 + 两条单位律）源于一个深刻的观察：

\begin{center}
\framebox{\textbf{伴随函子的复合 $GF$ 自动满足单子公理。}}
\end{center}

\subsection{历史与哲学线索}

\begin{enumerate}[nosep]
    \item \textbf{Roger Godement（1958）} 在研究层上同调时，首次形式化了“标准构造”（standard construction），后来被称为单子；
    \item \textbf{Peter Huber（1961）} 观察到，任意一对伴随函子都产生这样的结构；
    \item \textbf{Samuel Eilenberg \& John Moore（1965）} 系统化：给定单子 $\mathbb{T}$，构造其代数范畴（Eilenberg-Moore 范畴），并用 Kleisli 范畴给出每个单子都可表示为伴随函子的复合；
    \item \textbf{Saunders Mac Lane（1971）} 在《Categories for the Working Mathematician》中确立了伴随-单子-代数三位一体的标准叙述。
\end{enumerate}

\subsection{泛代数的联系}

单子公理对代数结构的编码能力：
\begin{itemize}[nosep]
    \item \textbf{群单子}：$T(X) = $ 由 $X$ 生成的自由群，$\eta_X$ 将生成元嵌入，$\mu_X$ 将“词中的词”展平为词（结合律 = 展平的不同顺序给出相同结果）；
    \item \textbf{Maybe 单子}（计算机科学）：$T(X) = X + 1$（添加异常值），单位律确保异常传播的正确性；
    \item \textbf{State 单子}（编程）：$T(X) = (X \times S)^S$，结合律确保状态传递幂等性。
\end{itemize}

所有这些例子背后，$\eta$（单位 = \texttt{return}）和 $\mu$（乘法 = \texttt{join}）的来源正是伴随函子的结构投射。

\subsection{伴随-单子-余单子的三角关系}

\[
\begin{tikzcd}[column sep=5em, row sep=4em]
& \text{伴随} \; F \dashvadj G \arrow[dl,Rightarrow] \arrow[dr,Rightarrow] & \\
\text{单子} \; GF \arrow[rr,Leftrightarrow,"\text{Kleisli/Eilenberg-Moore}"] & & \text{余单子} \; FG
\end{tikzcd}
\]

每个单子都由（至少）两个伴随函子对生成：Kleisli 伴随 $F_K \dashvadj G_K$ 和 Eilenberg-Moore 伴随 $F_{EM} \dashvadj G_{EM}$，前者是初始（initial）的，后者是终端（terminal）的。这是范畴论的又一深刻定理，将在后续笔记中展开。

\newpage

% ==================================================================
\section{总结}
% ==================================================================

\subsection{核心要点}

\begin{enumerate}[nosep]
    \item \textbf{伴随的两种定义等价}：Hom-集自然同构 $\cong$ 单位--余单位 + 三角形恒等式。前者便于计算，后者便于构建全局结构。
    \item \textbf{伴随自然诱导单子}：$T = GF$，$\eta$ 直接复用为单子单位，$\mu = G\varepsilon F$ 为单子乘法。结合律源于 $\varepsilon$ 的自然性，单位律直接来自三角形恒等式。
    \item \textbf{伴随自然诱导余单子}：$S = FG$，$\varepsilon$ 直接复用为余单子余单位，$\delta = F\eta G$ 为余单子余乘法。公理同样源自伴随条件。
    \item \textbf{单子/余单子不是凭空定义的}：它们的三条公理是伴随函子的三角形恒等式和自然性的直接投影。这就是单子的真正动机。
\end{enumerate}

\subsection{后续学习路线}

掌握本讲内容后，后续笔记将覆盖：
\begin{itemize}[nosep]
    \item Kleisli 范畴与 Eilenberg-Moore 代数：每个单子如何还原为伴随；
    \item 自由-遗忘伴随的具体例子（群、环、模等）；
    \item 单子变换（Monad Transformers）与复合；
    \item 编程中的单子（Haskell 的 Monad 类型类，do-notation 的范畴语义）。
\end{itemize}

\vspace{1em}
\begin{center}
\rule{0.5\textwidth}{0.5pt}\\[0.5em]
\textbf{本节完}
\end{center}

\end{document}